\chapter{Register of Claims and Their Status}
\label{app:register}

The book mixes three kinds of statement: results proved or cited, definitions and stipulated examples, and hypotheses about perception and craft. The register below lists the hypotheses, with the part of each that is mathematics, the part that is empirical, and what would count against it. Every counting and numerical claim of the worked examples (the counts of edits, orderings, and realizations, the occupation tables, the singular values, and the remainder bounds) is reproduced by the script \texttt{verify.py} distributed with the source, computed independently of the text; a disagreement between script and text indicates an error in one of them. Numerical values in worked examples (the forged-letter scene, the response model, the sound levels) are stipulated for illustration and are not measurements.

\begin{center}
\small
\begin{tabular}{p{2.6cm}p{3.2cm}p{3.5cm}p{3.6cm}}
\toprule
principle & mathematical part & empirical part & would count against it\\
\midrule
Geodesic principle (\Cref{pr:geodesic}) & Euler--Lagrange equation of \Cref{prop:EL} & that natural camera motion approximately extremizes such an action & natural paths failing the equation for every metric in the family\\[4pt]
Obstruction principle (\Cref{pr:obstruction}) & equivalence of consistency with a vanishing class, once corrections lie in an abelian sheaf & that a film's continuity constraints can be modeled that way & a film whose constraints do not fit, or a repair by local edits of a nonzero class\\[4pt]
Authority principle (\Cref{pr:authority}) & none beyond the layer definitions & that frame-only systems fail because they lack a history layer & failures of such systems unrelated to missing history, or success without the layer\\[4pt]
Shared-segment principle (\Cref{pr:shared}) & typing in an intersection of contexts & that authors can state the contexts of a branching film & shared segments that are well typed on some branches and still acceptable on all\\[4pt]
Conditional dimension (\Cref{pr:conditional}) & first-order effective dimension (\Cref{prop:effdim}) and its remainder (\Cref{prop:taylor}) & that the number of independent variations depends on scene, viewer, task, tolerance, and budget & estimates that are stable across scenes and populations to a degree justifying a single figure\\
\bottomrule
\end{tabular}
\end{center}

Two further conjectures are stated in the text without being named as principles. The conjecture of Chapter~\ref{ch:equivalence} that $d_{\mathrm{eff}}$ is small relative to the number of controls near typical production settings is empirical, and a slowly decaying singular spectrum would count against it. The claim of Chapter~\ref{ch:contracts} that listener preference is not determined by reduced occupation, preserved intelligibility, and preserved cues is also empirical, and is stated as a warning against inferring the fourth from the first three.
